\chapter{Discourse Structure and Spoken Usage}
\label{ch:discourse}\label{ch:spoken}\label{ch:semantics}

% ============================================================
% This chapter brings together three domains that belong to a single
% descriptive continuum: how Iridian organises information across clauses
% and texts; how register and social context shape grammatical choices;
% and how the spoken and colloquial variety departs from the written
% standard. The three domains interact: aspect choice as a narrative
% device is also a discourse-register choice; many spoken morphemes are
% formally integrated into the same inflectional machinery that the rest
% of the grammar describes; and politeness strategies depend on the same
% anchoring and topic mechanisms analysed in earlier chapters.
% ============================================================

\section{Introduction}\label{sec:discourse-intro}

Iridian maintains reference across clauses through overt nominals,
 demonstratives, reflexives, and recoverable null arguments. Aspect and voice
 determine the grammatical pivot of each predicate, while discourse determines
 which participant remains the current subject matter. The two often
 coincide, although linked agent topics and external resumptive topics permit
 their separation (Chapter~\ref{ch:infstr}). Register and speaker stance
 further condition the choice of particles and the use of colloquial forms.

%
% ================================================================
%  PART I: DISCOURSE STRUCTURE
% ================================================================
%

\section{Topic continuity and topic shift}\label{sec:topic-continuity}


\subsection{Reference management without morphological cross-reference}
\label{sec:ref-management}

Iridian manages reference tracking without any pronominal clitic series on the
verb or noun. The devices available are: (a)~overt demonstratives, which encode
the proximate/obviative distinction and serve as the primary third-person
referential series; (b)~overt first- and second-person pronouns, employed when
required by contrast, emphasis, disambiguation, or low discourse recoverability;
(c)~the reflexive anaphor, which forces local coreference within the clausal
domain; and (d)~zero anaphora, the unmarked strategy in contexts of stable topic
continuity, where the intended argument is recoverable
without any overt expression.

Iridian has no clitic or agreement layer that morphologically
cross-references verbal arguments: argument omission depends on discourse
recoverability within the licensing conditions of each construction. Directive
applied participants and causees remain obligatorily overt
(§\,\ref{sec:mood-directive-valency}), even when their referents are
recoverable. The verb supplies no right-edge person morphology to
disambiguate an otherwise licensed null argument.

Lexical activity readings and service causatives have different argument
structures from ordinary clauses with unexpressed referential arguments.
The indicative NOM-domain activity uses of \ird{pjašt-} and \ird{kař-} license no
separate patient (§\,\ref{sec:nom-activity}), while the restricted service
construction leaves its lower performer existentially understood
(§\,\ref{sec:service-causative}). Neither interpretation depends on
recovering that participant as a continuing discourse topic.

An overtly reintroduced topic may receive a modest pitch reset when it
 opens a new discourse segment, without thereby contrasting alternatives.
 Contrastive topicality instead relates parallel subquestions and permits
 the marked rising or fall-rise patterns described in
 §\,\ref{sec:tc-contrastive-topic}. A continuing topic can remain
 deaccented or unexpressed; neither option requires an audible boundary
 before the comment (§\,\ref{sec:phrasal-prominence}).

Implicit nominal possession has its own licensing conditions
(§\,\ref{sec:nm-implicit-poss}). In \ird{Čás konca} \trsl{head hurts},
the noun remains overt while its relational possessor is contextually
resolved. An accessible topic or utterance perspective can favour a referent,
but cannot override an overt possessor or an independently licensed binder.
The grammatical pivot and converbial pivot continuity do not independently
fix possessor identity; nominal anchoring does not encode speaker ownership.

\section{Anchoring in discourse context}\label{sec:discourse-definiteness}

% The anchoring suffix -to on a noun signals that the referent is
% identifiable to the hearer. In discourse, first mention is typically
% bare (no anchoring marker); subsequent mention may use the marker,
% a demonstrative, or drop the noun entirely (zero anaphora).
% The interaction of anchoring, topic, and absolutive alignment
% is the key thread here.

\subsection{The neutral discourse centre in direct-case position}
\label{sec:disc-def}

Shared anchoring, expressed by suffixal \ird{-to}/\ird{-te} and their
 allomorphs or by a class-specific cumulative ending, maintains an identifiable
 nominal referent, whether that referent is a topic, a backgrounded participant, or
 part of the asserted information. Direct-case patients and transitive-case
 agents have the same shared-accessibility contrast; their case does not add
 contrastive or resumptive meaning to the anchor
 (§\,\ref{sec:nm-def-nondirect}). A linked topic can independently promote
 an anchored agent to the subject matter of a perfective clause, while the
 patient retains direct case (§\,\ref{sec:tc-marked-topics}).

A continuing aboutness topic supplies a strong contextual antecedent for a
 recoverable null argument. Where no such topic has been independently
 established, the preceding direct-case pivot provides the neutral fallback.
 These are defeasible discourse preferences: the predicate's interpretation,
 competing referents, and overt reintroduction constrain reference. They do
 not replace the same-pivot relation in converb chains
 (§\,\ref{sec:switch-reference}).

\section{Aspect choice as narrative strategy}\label{sec:aspect-narrative}

Perfective and retrospective clauses commonly advance a sequence of bounded
 events, whereas imperfective clauses can describe ongoing activities or
 background states. Their different case frames also change the default
 pivot candidate for discourse continuation. That association does not force
 a change of topic: a perfective agent can remain the subject matter of an
 independently topicalised predication, and an imperfective description can
 introduce a new referent (§\,\ref{sec:tc-marked-topics}).

The antipassive supports agent continuity by retaining the agent as pivot
while removing the patient from core argument structure. It presents an
event chiefly as part of the agent's activity, allowing that agent to remain
topical without establishing or tracking a core patient. It is especially natural when
the patient is generic, incorporated into the event type, inferable from
context, or simply of no further narrative consequence; for the morphological
details see Chapter~\ref{ch:verbs}.

The causative introduces a causer and presents the event as one of
inducement, delegated action, permission, or coercion. In narrative, this
additional participant can foreground authority, distribute responsibility,
or distinguish the discourse topic from the participant who directly
performs the action.
Neither perfective nor retrospective aspect determines whether the causation
is coercive, permissive, direct, or mediated; these interpretations depend
on the predicate and context (§\,\ref{sec:aspect-caus}). Although the added
causer supplies another potential centre for discourse continuation, its
selection as topic remains independent of the causative morphology.

\subsection{Valency and the choice of continuing participant}
\label{sec:disc-valency-continuity}

A perfective transitive, an antipassive, and an applicative offer different
participants as their direct-case pivots, although an independently established
topic may differ from each. In the following constructed sequences, a
perfective buying clause is continued with a description of its patient, an
antipassive eating clause with a description of its actor, and an applicative
with a description of its beneficiary. Each continuation repeats the intended
participant overtly; the examples illustrate the availability of these
continuations without establishing a preference between competing null
antecedents.

\pex\label{ex:disc-valency-sequences}
\a
\begingl
\gla Travó Marku koupnik. Travó jeca.//
\glb bread.\Dir{}.\Def{} Marek.\Trs{} buy-\Trs{}-\Pf{} bread.\Dir{}.\Def{} \Exst{}-\Cont{}//
\glft \trsl{Marek bought the bread. The bread is there.}//
\endgl
\a
\begingl
\gla Marek pjášček. Marek zdrávime.//
\glb Marek.\Dir{} eat-\Antip{}-\Pf{} Marek.\Dir{} sleep-\Prog{}//
\glft \trsl{Marek ate. Marek is sleeping.}//
\endgl
\a
\begingl
\gla Byl Marku travě pjaštébik. Byl zdrávime.//
\glb child.\Dir{} Marek.\Trs{} bread.\Ind{} eat-\Appl{}-\Pf{} child.\Dir{} sleep-\Prog{}//
\glft \trsl{Marek ate bread for the child's benefit. The child is sleeping.}//
\endgl
\xe

\noindent The first sequence retains a patient pivot, the second an agent
pivot, and the third an applied beneficiary. The antipassive excludes an overt
patient from its finite clause; the applicative retains the theme in indirect
case while selecting the beneficiary. In each sequence the repeated
participant can serve as the continuing subject matter in a suitable context.
A null continuation remains subject to discourse recoverability and the selectional
requirements of its own predicate (§\,\ref{sec:ref-management}); applicative
morphology alone neither establishes a topic nor excludes reference to the
agent or theme.

\subsection{Mini-texts: continuity and re-centering}

The following constructed sequences illustrate topic continuity where the
 topic coincides with the pivot, and a subsequent shift supported by overt
 reintroduction.

\pex
\a
\begingl
\gla Marek travě pjaštíme.//
\glb Marek.\Dir{} bread.\Ind{} eat-\Prog{}//
\glft \trsl{Marek is eating bread.}//
\endgl
\a
\begingl
\gla Zemácime.//
\glb run-\Prog{}//
\glft \trsl{[He] is running.} \lingcontext{agent continuity in the nominative domain}//
\endgl
\xe

\noindent Here the null topic of the second clause is naturally interpreted as
Marek, who is both the current aboutness topic and the direct-case pivot.
The predicate \ird{zemac-} also favours an animate participant, so the
inanimate bread is not an equivalent antecedent competitor.

\pex
\a
\begingl
\gla Tak tóm Marku koupnik.//
\glb \Dem{}.\Prox{}.\Dir{} book.\Dir{} Marek.\Trs{} buy-\Trs{}-\Pf{}//
\glft \trsl{This book, Marek bought.}//
\endgl
\a
\begingl
\gla By koupnik?//
\glb \textsc{q} buy-\Trs{}-\Pf{}//
\glft \trsl{Did [he] buy [it]?}
 \lingcontext{confirmation question continuing the same buyer and book}//
\endgl
\xe

\noindent In an exchange about the book, the confirmation question continues
 the patient topic while leaving the buyer unexpressed as well. The same
 perfective case frame holds in both clauses, and the discourse recovers each
 participant in the role selected by the verb. Since both referents remain
 accessible, the sequence illustrates continuity without isolating a special
 antecedent preference for the pivot. Overt reintroduction can instead make
 an agent the subject matter of the next clause, as in the following sequence.

\pex
\a
\begingl
\gla Travó Marku pjaštník.//
\glb bread.\Dir{}.\Def{} Marek.\Trs{} eat-\textsc{trs}-\Pf{}//
\glft \trsl{Marek ate the bread.}//
\endgl
\a
\begingl
\gla Marek travě pjaštíme.//
\glb Marek.\Dir{} bread.\Ind{} eat-\Prog{}//
\glft \trsl{Marek is eating bread.} \lingcontext{re-centering from patient-topic narration to agent-topic description}//
\endgl
\xe

\noindent The first clause can describe what happened to the bread, while
 the second explicitly reintroduces Marek as the participant whose current
 activity is described. The aspectual alternation and overt nominal jointly
 support that shift; neither variable has been held constant.

The established agent-topic clause \ird{Markuto tak tóm koupnik}
 (\ref{ex:tc-marked-agent-top}) separates the relations more directly:
 Marek is the aboutness topic and the book the grammatical pivot. A later null
 agent can be recovered as Marek when the discourse continues his activities,
 although a null patient in a matching buying clause can instead recover the
 book. The intended referent depends on its role in that clause, rather than
 on a mechanical instruction to copy the last direct-case expression.

The relative preference for a non-pivot topic and a competing pivot remains
 unresolved when both referents are equally plausible antecedents. The
 preceding sequences differ in animacy, predicate choice, or overt expression;
 they establish recoverability and overt shift, while leaving that competition
 under fixed aspect, case, and anchoring undetermined. The distinction from converb continuity
 remains grammatical: its same-pivot dependency is defined over the two
 predicates' pivots rather than over whichever participant the discourse is
 about (§\,\ref{sec:switch-reference}).

\subsection{Controlled continuations with competing animate referents}
\label{sec:disc-controlled-continuations}

Shared anchoring can be held constant when a transitive agent and a
 direct-case patient compete as possible antecedents. In the following
 constructed pair, both Marek and Janek are already identifiable; the
 contextual question selects either what happened to Janek or what Marek did.
 The case frame and perfective predicate remain identical, while the
 agent-topic order follows the linked construction in
 §\,\ref{sec:tc-marked-topics}.

\pex\label{ex:disc-animate-competition}
\a\label{ex:disc-patient-centre}%
\begingl
\gla Janko Markuto védnik.//
\glb Janek.\Dir{}.\Def{} Marek.\Trs{}.\Def{} see-\Trs{}-\Pf{}//
\glft \trsl{Marek saw Janek.}
 \lingcontext{Janek-topic context: what happened to Janek?}//
\endgl
\a\label{ex:disc-agent-centre}%
\begingl
\gla Markuto Janko védnik.//
\glb Marek.\Trs{}.\Def{} Janek.\Dir{}.\Def{} see-\Trs{}-\Pf{}//
\glft \trsl{As for Marek, he saw Janek.}
 \lingcontext{Marek-topic context: what did Marek do?}//
\endgl
\xe

\noindent Either clause can be followed by the same intransitive continuation:

\ex\label{ex:disc-ambiguous-continuation}%
\begingl
\gla Zdrávime.//
\glb sleep-\Prog{}//
\glft \trsl{[He] is sleeping.}//
\endgl
\xe

\noindent Both antecedents are animate and compatible with sleeping, so that
 lexical selection no longer resolves the reference. The topic-continuity
 analysis predicts Janek as the continuing subject matter after
 (\ref{ex:disc-patient-centre}) and Marek after
 (\ref{ex:disc-agent-centre}), whereas a preference determined solely by the
 preceding direct-case pivot would favour Janek after both. These are
 contrasting predictions of the analyses, rather than established judgements
 about the unexpressed sleeper. The strength of either preference, and the
 contextual conditions under which the competing referent remains available,
 are unresolved. Overt \ird{Marko} or \ird{Janko} before \ird{zdrávime}
 specifies the sleeper without requiring that preference to settle the choice.
 The finite continuation has its own progressive case frame; the perfective
 case and anchoring of both antecedents are controlled in the preceding pair.

With Marek indexed $i$ and Janek indexed $j$, the two readings of the
continuation are accordingly $\varnothing_i$ \ird{zdrávime} and
$\varnothing_j$ \ird{zdrávime}. Changing the preceding topic holds these
candidate antecedents constant while testing a discourse preference. It does
not create a local binding relation between the two finite clauses. In
particular, neither candidate licenses a reflexive in the continuation merely
by being the preceding pivot or topic: the local antecedent requirement for
\ird{svem} is a separate grammatical dependency
(§\,\ref{sec:nm-pron-refl}). A recoverable silent argument and a locally
bound anaphor thus require different evidence, even when their intended
referents coincide.

An overt topic can also identify the return to an earlier subject matter
 after an intervening predication about someone else. The following sequence
 holds the progressive aspect and shared nominal anchoring constant:

\pex\label{ex:disc-reactivation}
\a
\begingl
\gla Marko domů chódime.//
\glb Marek.\Dir{}.\Def{} home go-\Prog{}//
\glft \trsl{Marek is going home.} \lingcontext{Marek is the current topic}//
\endgl
\a
\begingl
\gla Janko zdrávime.//
\glb Janek.\Dir{}.\Def{} sleep-\Prog{}//
\glft \trsl{Janek is sleeping.} \lingcontext{shift to Janek}//
\endgl
\a
\begingl
\gla Marko domů so chódime.//
\glb Marek.\Dir{}.\Def{} home still go-\Prog{}//
\glft \trsl{Marek is still going home.} \lingcontext{return to Marek after the interruption}//
\endgl
\xe

\noindent Repetition of \ird{Marko} identifies the returning referent, while
 \ird{so} presents the journey as continuing. Shared anchoring is compatible
 with both the shift to Janek and the return to Marek; it does not distinguish
 those discourse transitions. Removing the final \ird{Marko} would leave
 reference dependent on the competing antecedents and the activity described.
 The overt sequence therefore illustrates reactivation without establishing
 that overt reintroduction is obligatory after every interruption, or that
 the suffix itself encodes reactivation.

\subsection{Nominal anchoring under a constant topic interpretation}
\label{sec:disc-anchoring-controlled}

A child can remain the subject matter while its nominal expression varies
 between bare, speaker-anchored, and shared-anchored forms. Under a discussion
 of a particular child's current activity, the following constructed clauses
 retain the same predicate, case, and topic--comment organisation:

\pex\label{ex:disc-anchoring-controlled}
\a
\begingl
\gla Byl zdrávime.//
\glb child.\Dir{} sleep-\Prog{}//
\glft \trsl{The child is sleeping.} \lingcontext{child-topic context; anchoring left unmarked}//
\endgl
\a
\begingl
\gla Si byl zdrávime.//
\glb \Sanch{} child.\Dir{} sleep-\Prog{}//
\glft \trsl{A certain child is sleeping.} \lingcontext{child-topic context; overt speaker selection}//
\endgl
\a
\begingl
\gla Bylo zdrávime.//
\glb child.\Dir{}.\Def{} sleep-\Prog{}//
\glft \trsl{The child is sleeping.} \lingcontext{child-topic context; shared-discourse retrieval}//
\endgl
\xe

\noindent The bare expression leaves the source of identification unmarked,
 \ird{si byl} overtly presents speaker selection, and \ird{bylo} invokes
 shared accessibility. The speaker-anchored form does not assert that the
 addressee is unable to identify the child. A context in which both
 interlocutors can identify the child permits the anchoring contrast to be
 considered without changing which individual is under discussion
 (§\,\ref{sec:nm-spec}); the choice of anchor changes the conventional
 presentation of that referent rather than assigning its topic status.
 A subsequent \ird{Zdrávime} can continue the child under each introduction
 when reference remains recoverable. No difference in the strength of that
 continuity preference is established by these forms alone.

\section{Discourse particles}\label{sec:discourse-particles}\label{sec:discourse-word-order}


\subsection{Clause-final discourse particles}\label{sec:disc-final-particles}

Iridian permits a small class of clause-final discourse particles, placed after
the final prosodic constituent of the clause, most typically after the finite
verbal complex. These particles do not belong to the ordinary verbal
inflectional template and do not affect argument structure, alignment, or
truth-conditional content. Their function is instead pragmatic: they mark the
speaker's stance toward the proposition and toward the addressee at the level of
the utterance as a whole.

These particles stand at the utterance edge, outside the verbal stress domain
and the preverbal ordering template of §\,\ref{sec:particle-combinations}.
In weak uses they may lean on the final prosodic word of the preceding clause;
emphatic uses permit independent prominence. Neither attachment nor emphasis
relocates lexical stress within that host, and a weakly attached particle
contributes no syllable to the host's subsidiary footing
(§\,\ref{sec:secondary-stress}). A weak particle within the same
intonational phrase may carry its final boundary target without receiving
a nuclear pitch accent. In neutral predicate-focus delivery, the principal
accent remains on the preceding lexical host; with earlier narrow focus,
it can remain on that focused constituent. In either case the particle
supplies the final syllable on which L\% or H\% is realised. An emphatic particle can instead form an accentual
phrase with its own pitch accent; where it constitutes a separate
intonational phrase, it also has its own nuclear contour. Neither emphasis
nor utterance-edge position alone requires an intervening pause
(§§\,\ref{sec:prosodic-domains} and~\ref{sec:intonation}). No single
obligatory melody is assigned to each particle. Semantically, they express such
values as appeal to shared knowledge, assertive presentation, corrective
insistence, or emotive intensification.

The principal clause-final particles are the following.

\noindent\textit{\ird{len}}\quad  The particle \ird{len} is confirmatory and shared-ground
oriented. It presents the proposition as accessible, recognisable, or available
for the addressee's ratification, with values such as `you know', `as we both
know', or `is that not so'. It is the least forceful member of the set and is
well suited to mild prompting, reassurance, and alignment of speaker and
addressee.

\noindent\textit{\ird{prav}}\quad  The particle \ird{prav} is assertive and presentational.
It marks the utterance as one the speaker puts forward with particular
commitment, roughly `indeed', `truly', or `mark this'. As a discourse particle
rather than a lexical predicate, it strengthens the assertion without taking an
argument of its own. In
clause-final position it signals strengthened assertion without necessarily
implying contrast.

\noindent\textit{\ird{však}}\quad  The particle \ird{však} is corrective or contrastive. It
indicates that the utterance stands against an assumption, implication, or line
of expectation in the discourse, with values such as `after all', `but in fact',
`mind you', or `contrary to that'. It is especially natural in rebuttal,
clarification, argumentative discourse, and narrative correction.

\noindent\textit{\ird{blěc}}\quad  The particle \ird{blěc} is a marked colloquial and
low-register emotive particle. Unlike \ird{len}, \ird{prav}, and \ird{však}, it
does not belong to the neutral core discourse system, but adds rough speaker
attitude such as irritation, incredulity, impatient insistence, or coarse
emphasis. Its contribution is primarily expressive and register-sensitive. In
formal or careful style it is normally avoided; in colloquial speech it may
reinforce commands, corrections, or exasperated assertions.

\subsection{Co-occurrence and interactional scope}
\label{sec:final-particle-combinations}

The utterance-final distribution does not specify whether two of these
particles can co-occur, their relative order, or their availability with
each preverbal stance particle and clause type. The overlap between
affirmative \ird{mó} and assertive \ird{prav}, or between
discourse-reorienting \ird{jas} and corrective \ird{však}, establishes
neither complementary distribution nor a conventional reinforcing
combination. An appeal for the addressee's ratification, an expression of
speaker commitment, and a correction of an earlier assumption have
different interactional targets, even when they concern the same
proposition. The individual meanings above do not determine whether those
targets can be jointly expressed by an Iridian particle sequence.

Prosodic attachment likewise concerns the final particle's local host,
whereas interactional scope concerns the contribution it qualifies.
The parenthetical hedge \ird{žma} and the transition connective
\ird{nu} have their own distributions
(§§\,\ref{sec:spoken-zema}, \ref{sec:spoken-nu}); their
utterance-edge placement supplies no general ordering rule with the final
particles. Nor is a specific final rise, fall, or intervening pause assigned
to an unspecified combination by the independent prominence available to
its members.

Japanese illustrates the independence of these restrictions. The order
\textit{ka} $<$ \textit{yo} $<$ \textit{ne} permits combinations whose interpretation depends
on clause type and intonation: \textit{ka yo} has a rhetorical reading
with a final fall and excludes polite verbal forms in the treatment of
McCready and Davis, although each particle separately permits those forms.
Their discussion of \textit{ka yo ne} further distinguishes a possible
particle cluster from a sequence followed by an independently phrased
confirmation request \cite[§5, manuscript pp.~21--25]{mccreadydavis2019particles}.
These distinctions supply comparative diagnostics for separating
co-occurrence, register, phrasing, and interactional target; they establish
no corresponding restriction or combination in Iridian.

%
% ================================================================
%  PART II: REGISTER AND PRAGMATICS
% ================================================================
%

\section{Spoken Iridian: an overview}\label{sec:spoken-overview}

This section sketches a small part of the spoken and urban colloquial register
of present-day Iridian. The written standard remains the reference variety of
this grammar, but it would be misleading to treat everyday spoken usage as
merely a noisier version of that standard. In ordinary conversation, speakers
make frequent use of reduced connectives, clause-peripheral hedges, semantically
bleached nouns, and other high-frequency colloquial forms whose contribution is
partly lexical and partly interactional.

Colloquial \ird{nu}, beside sequential \ird{nesu}, and the intentional use
of the purposive converb illustrate the interaction between lexical reduction
and discourse organisation (§§\,\ref{sec:seq-copula},
\ref{sec:purp-intentional}). The morphemes described below similarly combine
ordinary inflection with lexical or interactional meanings restricted to spoken
registers.

\section{Common colloquial morphemes}\label{sec:spoken-morphemes}

\subsection{The connective \ird{nu}}\label{sec:spoken-nu}

Colloquial \ird{nu} relates an utterance to the preceding discourse through
continuation, progression, or transition. Readings such as \trsl{and then},
\trsl{so}, and \trsl{well then} arise from the relation between the discourse
contributions: the connective itself does not require chronological succession
or encode a strong inference. A consequential interpretation is available
where the preceding content supplies grounds for the next move, while a change
of subject matter can give a weaker transition reading.

The colloquial reduction from copular converb \ird{nesu}
(§\,\ref{sec:seq-copula}) motivates the form's connection with clause
linkage. Synchronically, discourse \ird{nu} neither attaches to a dependent
predicate nor requires an unexpressed copular event; it stands outside the
eight-member productive converb paradigm.

The following constructed sequence contains two independent finite
sentences, with \ird{nu} introducing the second:
\ex\label{ex:nu-independent-continuation}%
\begingl
\gla Marek pjášček. Nu, Janek znóvime.//
\glb Marek.\Dir{} eat-\Antip{}-\Pf{} \textsc{disc} Janek.\Dir{} study-\Prog{}//
\glft \trsl{Marek ate. And Janek is studying.}
\lingcontext{a continuation concerning Janek, without a necessary causal relation}//
\endgl
\xe
Both sentences retain the morphology of independent predications; neither
bears finite dependent \ird{-e}. A preceding full utterance boundary is
compatible with this use of \ird{nu}, although the connective alone specifies
neither a pause duration nor a pitch-reset magnitude. Its discourse-edge
position differs from the position of preverbal stance particles
(Chapter~\ref{ch:function-words}).

Consequential \ird{nadě} already presents a proposition as the upshot of facts
in view, while \ird{hladě} marks an evidence-grounded inference
(§\,\ref{sec:class3-particles}). These particles qualify the following
predication within the preverbal field; \ird{nu} links the discourse move to
what precedes it. An English translation with `so' can therefore render either
a weak transition or an explicit consequential stance, depending on the
Iridian construction. Their distinct distributions do not establish that
all combinations of \ird{nu} and stance particles are licensed.

\subsection{Continuation, inference, and justification}
\label{sec:spoken-justification-examples}

An independent continuation can relate events in their chronological order
or present one proposition as evidence for another. In the lexical-slot
schemata below (§\,\ref{sec:lexical-slot-schemata}), \ird{nu} permits a
contextual consequence, whereas \ird{hladě} explicitly presents an inference
from visible evidence:

\pex\label{ex:slot-discourse-relations}
\a
\begingl
\gla Marek ščének. Nu, Janek \textnormal{[[depart]]}\textsubscript{\Pf{}}.//
\glb Marek.\Dir{} arrive-\Pf{} \textsc{disc} Janek.\Dir{} depart-\Pf{}//
\glft \trsl{Marek arrived. Then Janek departed.}
\lingcontext{arrival may motivate departure, but the continuation does not require that cause}//
\endgl
\a
\begingl
\gla \textnormal{[[road]]}\textsubscript{\Dir{}} \textnormal{[[be-wet]]}\textsubscript{\Cont{}}. \textnormal{[[rain]]}\textsubscript{\Dir{}} hladě \textnormal{[[fall]]}\textsubscript{\Pf{}}.//
\glb road.\Dir{} be.wet-\Cont{} rain.\Dir{} infer fall-\Pf{}//
\glft \trsl{The road is wet. It must have rained.}
\lingcontext{wetness supplies evidence for an earlier rain event}//
\endgl
\a
\begingl
\gla \textnormal{[[rain]]}\textsubscript{\Dir{}} \textnormal{[[fall]]}\textsubscript{\Prog{}}. \textnormal{[[umbrella]]}\textsubscript{\Dir{}} \textnormal{[[take]]}\textsubscript{\textsc{imp.ipfv}}.//
\glb rain.\Dir{} fall-\Prog{} umbrella.\Dir{} take-\Imp{}.\textsc{ipfv}//
\glft \trsl{It is raining. Take an umbrella.}
\lingcontext{rain motivates the instruction; its addressee is recoverable as the omitted second-person actor}//
\endgl
\a
\begingl
\gla \textnormal{[[door]]}\textsubscript{\Dir{}} Marku \textnormal{[[open]]}\textsubscript{\textsc{trs.pf}}. Janek nadě \textnormal{[[enter]]}\textsubscript{\Ctpl{}}.//
\glb door.\Dir{} Marek.\Trs{} open-\Trs{}-\Pf{} Janek.\Dir{} conseq enter-\Ctpl{}//
\glft \trsl{Marek opened the door. Consequently, Janek will enter.}
\lingcontext{entry is presented as the upshot of the preceding fact}//
\endgl
\xe

The evidence in (b) supports a conclusion about an earlier event; it
does not make road wetness a cause of that event. In (c), the first assertion
instead justifies the following instruction. The imperative cell is imperfective and theme-oriented, with a direct
patient and a recoverable transitive-case actor under directive Mood
(§\,\ref{sec:mood-directive-case}); its inflection is internal to that
independent sentence. The relation concerns
the reason for issuing the instruction, rather than causation of taking
within the represented situation. In (d), the completed opening has a direct
patient and transitive-case actor, while \ird{nadě} presents the next
proposition as its contextual upshot. That consequential stance is more
explicit than the weak transition supplied by \ird{nu} in (a).

Each sequence contains independent finite sentences, without dependent
\ird{-e}. Neither independently scoped directives inside a finite
conjunctive complex nor a propositional-anaphoric exponent follow from these
discourse relations.

\section{Propositional anaphora and connective expressions}
\label{sec:propositional-anaphora}

An event, proposition, or discourse act may constitute a different anaphoric
target from a participant mentioned within a clause. Demonstrative systems
can distinguish participant tracking from reference to a preceding clause or
proposition; demonstrative-based conjunctive adverbs provide a further
cross-linguistic connection between anaphora and clause linkage
\cite[306,321--325]{diesselbreunesse2020}. In Iridian, the established reduced
series \ird{i}/\ird{ji}/\ird{ni} resumes highly accessible nominal referents
(§\,\ref{sec:nm-dem-reduced}). That distribution does not yet establish
reference to the proposition expressed by an entire preceding utterance.

A propositional-anaphoric extension is specified here as a constructional
proposal, without a canonical lexical exponent. Its antecedent would be an
accessible proposition, rather than a participant selected from that
proposition. The construction would have to distinguish this target from an
event token and from the preceding speech act, and supply a predicate or
relational expression that licenses a proposition-valued argument. No
member of \ird{i}/\ird{ji}/\ird{ni} is assigned that function by default.

Three English schemas delimit the proposed connective family:
\begin{table}[ht]
\footnotesize
\caption{Proposed anaphoric connective relations without specified Iridian exponents.}
\label{tab:proposed-anaphoric-connectives}
\begin{tabularx}{\textwidth}{l L L}
\toprule
Schema & Antecedent status & Relation to the following discourse \\
\midrule
`That being so' & a premise accepted for the current move & circumstance supporting a continuation or conclusion \\
`If so' & a proposition entertained conditionally & restriction on the alternatives in which the continuation is evaluated \\
`Even so' & an accepted premise & continuation defeating an expectation associated with that premise \\
\bottomrule
\end{tabularx}
\end{table}

The schemas distinguish premise acceptance, conditional restriction, and
counter-expectation; they are not translated Iridian expressions. The proposed
constructions require an independently licensed anaphor and source predicate
before any combination with circumstantial, conditional, or concessive
morphology can be established. The resulting phrase must also function as a
connective over independent discourse contributions, rather than as an ordinary
converb clause describing an event. An inferential expression additionally
requires a conventional consequence relation beyond the weak continuation
already supplied by
\ird{nu}; existing \ird{nadě} and \ird{hladě} remain available in their
preverbal constructions.

Korean illustrates several distinct mechanisms: Sohn distinguishes
clause-final connective endings, conjunctive adverbials, and unmarked
apposition \cite[291--292]{sohn2009}. These categories do not map directly
onto Iridian's contrast between non-finite converbs, dependent finite
linkage, and independent discourse linkage. Nor does a conjunctive adverbial
necessarily require a full utterance boundary: the Korean adverbials can
co-occur with connective endings inside one intonational unit
\cite[290]{sohn2009}. The proposed Iridian family therefore draws on an
anaphoric constructional pattern without importing Korean lexical forms,
finiteness, or prosodic distributions.

\section{\ird{Žema} and reduced \ird{žma}}\label{sec:spoken-zema}

Among colloquial noun-like items, \ird{žema} is especially revealing. In
present-day spoken Iridian it is best analysed as a semantically bleached
colloquial lexeme with three closely related functions. First, it may serve as
a placeholder noun meaning roughly \trsl{thingy}, \trsl{that thing}, or
\trsl{what-do-you-call-it}. Secondly, it may function as a dismissive or
evaluative noun meaning \trsl{nonsense}, \trsl{excuses}, or \trsl{all that
fuss}. Thirdly, it may appear at the edge of an utterance as a playful hedge,
roughly \trsl{not really}, \trsl{just kidding}, or \trsl{I am only saying it
lightly}.

The reduced variant \ird{žma} shares these meanings, while its distribution
is associated with rapid, playful, or dismissive casual speech. This fixed
colloquial reduction is lexically specified; it does not extend the productive
deletion of an unstressed stem-final-syllable \ird{e} described in §\,\ref{sec:morphophon}.

\subsection{Category and core values}

The basic grammatical point is that \ird{žema} behaves as a noun-like base in
ordinary Iridian syntax. In its placeholder use it fills a nominal slot without
contributing a stable lexical class of its own:

\pex
\a
\begingl
\gla Tak žema Janku koupnik.//
\glb this thingy Janek.\Trs{} buy-\Trs{}-\Pf{}//
\glft \trsl{Janek bought that thing / that whatchamacallit.}//
\endgl
\a
\begingl
\gla Tak žema.//
\glb this nonsense//
\glft \trsl{That is nonsense / that is just an excuse.}//
\endgl
\xe

In the first example, \ird{žema} is simply a placeholder noun occupying the
patient slot of an ordinary transitive clause. In the second, the same form is
interpreted dismissively: the speaker rejects what has just been said or done as
empty, unserious, or merely performative. The lexical content is therefore weak
in both uses, but the syntactic behaviour remains nominal.

\subsection{Plurality}

The nominal behaviour of the item becomes still clearer when it appears in
plural form. Following the regular Class~III \ird{-a} pattern,
\ird{žema} takes the plural direct
\ird{žémy}, contrasting with singular transitive \ird{žemy}, the plural
transitive \ird{žém} (zero-ending on the strong consonant stem), and the
plural indirect \ird{žémam}. In this plural configuration \ird{žema} denotes a loose aggregate:
\trsl{all those things}, \trsl{all that stuff}, or, in a more evaluative
context, \trsl{all that nonsense}.

\ex
\begingl
\gla Žémy Marku koupnik.//
\glb thingy.\Dir{}.\Pl{} Marek.\Trs{} buy-\textsc{trs}-\Pf{}//
\glft \trsl{Marek bought all that stuff.}//
\endgl
\xe

\subsection{Clause-peripheral hedge use}

The same colloquial base may also appear clause-peripherally, especially in the
reduced form \ird{žma}. In this use it no longer behaves like an ordinary noun
inside the clause, but as a parenthetical hedge by which the speaker playfully
withdraws, softens, or reframes what has just been said. Even so, it should not
be treated as a member of the formal particle hierarchy. Its contribution is
interactional and parenthetical rather than structural in the narrow sense.

\pex
\begingl
\gla Janek vy ščének, žma.//
\glb Janek.\Dir{} already arrive-\Pf{} hedge//
\glft \trsl{Janek has already arrived, not really / just kidding.}//
\endgl
\xe

This utterance-edge use follows naturally from the same semantic bleaching seen
in the nominal uses. The speaker introduces a proposition and then immediately
marks it as playful, non-committal, or not to be taken at full face value.
Synchronically, then, \ird{žema} and \ird{žma} form a small colloquial system:
an underspecified noun in the syntax, an evaluative noun in stance-marking
predication, and a parenthetical hedge in clause-peripheral discourse use.

\subsection{Verbal use and productive
recategorisation}\label{sec:spoken-zema-verbal}

The noun-like base also enters verbal morphology as a colloquial stem,
whose vague lexical content permits a range of contextual interpretations.

The colloquial verbal stem is \ird{žemat-}, assigned to thematic class~I
(thematic consonant \ird{-t}), with antipassive stem \ird{žemáč-}. Like
other class~I stems, \ird{žemat-} takes the valency and
aspect-alignment morphology described in Chapter~\ref{ch:verbs}, although its
lexical content remains as indeterminate as that of the nominal base.
Its quantity pattern is Class~A: the second root nucleus /a/ is the
productive target, so replacement antipassive selects \ird{žemáč-};
continuous \ird{žemata} retains G0.

In its actor-focused (antipassive) use, the stem describes the agent as engaging
in vague, performative, evasive, or unserious activity. The core meaning is most
readily glossed as \trsl{to mess around}, \trsl{to make excuses}, \trsl{to put
on a show}, or \trsl{to fuss about without achieving anything}, though the exact
nuance depends on context in the same way as the nominal base does. The
antipassive perfective \ird{žemáček} reports a completed episode of such
activity; the antipassive retrospective \ird{žemáčav} presents it with
continuing present relevance and a characteristically evaluative charge---the
speaker frames the prior activity as having amounted to nothing or as having
been mere performance; the continuous \ird{žemata} describes it as an ongoing
state or habitual disposition:

\pex
\a
\begingl
\gla Anka žemáček.//
\glb Anka.\Dir{} žemat-\Antip{}-\Pf{}//
\glft \trsl{Anka was just messing around / making excuses.}//
\endgl
\a
\begingl
\gla Anka žemáčav.//
\glb Anka.\Dir{} žemat-\Antip{}-\Ret{}//
\glft \trsl{Anka has been making nothing but excuses.}//
\endgl
\a
\begingl
\gla Anka žemata.//
\glb Anka.\Dir{} žemat-\Cont{}//
\glft \trsl{Anka is always fussing / is the type to make excuses.}//
\endgl
\xe

In its patient-focused (ergative) use, the stem \ird{žemat-} takes the
transitive marker \ird{-n-} and acquires a closely related but distinct reading:
\trsl{to handle carelessly}, \trsl{to give only a vague or makeshift treatment
to}, \trsl{to fudge}, or \trsl{to treat as a trivial matter}. In the ergative
aspects the agent appears in the transitive case and the patient in the
directional:

\pex
\a
\begingl
\gla Anka tak žemátnik.//
\glb Anka.\Trs{} that.\Dir{} žemat-\Trs{}-\Pf{}//
\glft \trsl{Anka fudged that / handled that carelessly.}//
\endgl
\a
\begingl
\gla Anka tak žematnáv.//
\glb Anka.\Trs{} that.\Dir{} žemat-\Trs{}-\Ret{}//
\glft \trsl{Anka has been treating that carelessly the whole time.}//
\endgl
\xe

Semantically, all verbal uses share a stable core even as contextual readings
vary widely. The predicate never names a concrete or specific action; it instead
attributes to the agent a \emph{stance}---vagueness, evasiveness, theatricality,
or dismissal---and that stance is precisely what the nominal base also
contributes. \ird{Žemaček} in one context means \trsl{she was just messing
around}; in another, \trsl{he was making excuses}; in a third, \trsl{they were
fussing over trivial extras}. What remains constant is the evaluative frame: the
speaker characterises the activity not by what was done but by how it was done,
and the \textit{how} is always some variant of unserious, unfocused, or merely
performative.

The thematic alternation \ird{žemat-}/\ird{žemáč-} and the aspect forms
illustrated above follow the ordinary class~I pattern. The construction's
colloquial character resides in the stem's vague, evaluative meanings and
register, while its inflection determines alignment and viewpoint in the
usual way.

\section{The \ird{ty}-supplement in spoken interaction}
\label{sec:spoken-ty}

The appositional particle \ird{ty} and the supplement construction it heads
are introduced in the nominal chapter (\S\,\ref{sec:nm-appos}). The present
section treats the specifically interactional properties of the
right-peripheral variant, which is substantially more productive in spoken
than in written language and which serves a distinct set of discourse functions
in everyday speech.

\subsection{Referent-first ordering as a spoken strategy}
\label{sec:spoken-ty-strategy}

Spoken language imposes real-time production constraints that written language
does not. A speaker who launches a noun phrase cannot always plan its full
modifier structure before the head noun must be uttered; equally, elaborating
descriptions may become relevant only as the discourse develops, after the
referent has already been introduced. The right-peripheral \ird{ty}-phrase
directly serves both situations. The speaker commits to a head noun,
establishes the referent, and then appends the clarification or specification
via \ird{ty} as a subsequent informational unit:

\pex
\a \ird{pres ty že védnikový i}
   \trsl{the man---that is, the one I saw}
\a \ird{nuž ty múžna i}
   \trsl{the city---the beautiful one}
\a \ird{Karel ty mlazka}
   \trsl{Karel---my brother, you know}
\xe

\noindent This ordering reflects the referential logic of spoken discourse: the
head noun anchors the referent in the common ground and the \ird{ty}-phrase
adds identifying or descriptive content that the hearer can use to resolve or
enrich the reference. The construction is particularly valuable when the
describing material is long or syntactically complex, since a heavy prenominal
modifier stack---a nominalised relative clause, a multi-word adjectival phrase,
a postpositional elaboration---requires planning the entire modifier before
uttering the head noun. Right-peripheral \ird{ty} makes it possible to deliver
the head noun immediately and defer the specifying material.

\subsection{Repair and afterthought uses}
\label{sec:spoken-ty-repair}

The most distinctively spoken use of the \ird{ty}-supplement is as a repair
and afterthought device. A speaker may introduce a referent, continue the
utterance, and then append a \ird{ty}-phrase retroactively to specify it more
precisely:

\pex
\ird{pres védnik, ty dumě na i}
   \trsl{I saw the man---the one at the house}
\xe

% Add how to use with Not for disambiguation

In such cases the \ird{ty}-phrase is structurally detached from the preceding
clause or \synhead{NP}. It is not embedded inside the \synhead{NP} it
elaborates; it is appended to the right edge of the utterance as a freestanding
supplement, anchored to the earlier referent anaphorically via \ird{i}. The
supplement begins a new prosodic unit, which may be realised with a short
pause before \ird{ty} or a pitch reset; a falling-then-rising transition and
a lower pitch level in the supplement are possible realisations of this
detachment. The construction prescribes neither a unique repair melody nor
silence at every boundary (§\,\ref{sec:boundary-phenomena}).

This repair-like use extends to cases where the \ird{ty}-phrase follows a
complete finite clause and reaches back to a referent introduced earlier in the
utterance:

\pex
\a \ird{Védnik, ty dumě na i}
   \trsl{I saw [him]---the one at the house}
\xe

Here \ird{ty} is not repairing an error but elaborating a
discourse referent that was already recoverable from context (via zero anaphora
in the main clause), offering the hearer additional identifying information as
an aside.

As described in \S\,\ref{sec:nm-appos}, the anaphoric pronoun \ird{i}
appears inside the \ird{ty}-phrase whenever the supplement's descriptive content
is non-nominal and requires a pronominal head to attach to. In spoken use this
is the normal case: spoken \ird{ty}-supplements typically contain adjectival or
relative-clause material, and \ird{i} serves as the pronominal host that allows
this material to form a well-headed nominal phrase while observing
head-finality. The resulting pattern, \ird{ty} \textsc{[modifier i]}, is highly
productive and has a characteristic phrasal contour: \ird{ty} receives mild
prominence, the modifier retains ordinary lexical stress, and \ird{i} is
typically reduced or deaccented at the right edge of the supplement. The
relative prominence of these words preserves each lexical word's initial
primary stress and ordinary subsidiary footing, while deaccentuation of
\ird{i} does not itself establish clitic status
(§\,\ref{sec:phrasal-prominence}).

Right-peripheral \ird{ty}-phrases are unmarked in spoken interaction. Their
frequency, their use in repair contexts, and their role as a strategy for
deferring heavy modifier material all reflect production-driven motivations that
are native to spontaneous speech. In formal and written language the
construction remains available but is marked: a right-peripheral \ird{ty}-phrase
in writing carries a deliberate elaborative implication, foregrounding the
supplement as a rhetorically distinct informational act rather than presenting
it as an integrated modifier. Heavy adnominal material---nominalised relative
clauses that would be unwieldy prenominally---can motivate the construction even
in formal prose without special effect; a short adjective or a simple nominal
going postnominal in writing, by contrast, is interpretable as emphatic or
contrastive. The contrast with the unmarked prenominal strategy is always
available to the writer as a pragmatic resource.

\section{Interjections}

Interjections stand outside ordinary clause structure and are not part of the
preverbal particle sequence described in Section~\ref{sec:adv-particles}.
